\section{Rappresentazione dei segnali multimediali}
\subsection{Rappresentazione di segnali audio}
\subsubsection*{Percezione del suono}
Il sistema uditivo umano percepisce i suoni effettuando un'analisi in frequenza del segnale audio che giunge all'orecchio.
Può essere modellato come un banco di filtri che scompongono il suono in più bande di frequenza, ciascuna delle quali
è associata ad una fibra nervosa che trasmette l'informazione al cervello. L'ascolto di un tono puro attiverà un solo
filtro che ecciterà la fibra nervosa corrispondente.

\subsubsection*{Soglia di percezione acustica}
La soglia di percezione acustica corrisponde al livello minimo di potenza sonora che può essere percepito dall'orecchio
umano. Varia in funzione della frequenza del suono e cambia da persona a persona, in base all'età. In generale la soglia
di percezione acustica è più bassa (ovvero l'orecchio umano è più sensibile) per le frequenze comprese tra 500 Hz e 5 kHz
che corrispondono alla fascia del parlato umano.

La soglia di pressione acustica non è sempre costante, ma in funzione dei suoni percepiti attraverso due fenomeni di mascheramento:
\begin{itemize}
	\item \textbf{mascheramento in frequenza}: la presenza di un suono forte in una certa banda di frequenza fa aumentare
	la soglia di percezione acustica dei suoni alle frequenze adiacenti, per cui un suono debole nella stessa banda di
	frequenza può non essere percepito
	\item \textbf{mascheramento temporale}: la modifica della curva di soglia di percezione acustica dovuta alla presenza
	di un suono forte in una certa banda di frequenza si propaga nel tempo per 2-5ms prima del suono forte e per 100-200ms
	dopo il suono forte, per cui un suono debole in una banda di frequenza può non essere percepito se è stato appena emesso
	un suono forte in una banda di frequenza vicina
\end{itemize}

\subsubsection*{Codifica percettiva}
La codifica percettiva sfrutta le caratteristiche del sistema uditivo umano e i fenomeni di mascheramento per ridurre
la quantità di dati necessaria a rappresentare un segnale audio senza introdurre distorsioni percepibili dall'orecchio
umano. Si può realizzare attraverso un'analisi tempo-frequenza che applica riquantizzazioni alle frequenze mascherate.

\subsubsection*{Digitalizzazione tramite Pulse Code Modulation (PCM)}
La codifica PCM consiste nella digitalizzazione di un segnale audio analogico attraverso le sole fasi essenziali di
campionamento, quantizzazione e inverse bit mapping. Non vengono applicate compressioni, per questo viene anche detta
codifica lossless. Gli unici parametri che influenzano la qualità del segnale digitale sono, quindi, la frequenza di
campionamento \(f_c\) e il numero di livelli di quantizzazione \(L\).

\subsubsection*{Digitalizzazione di un segnale vocale}
Per digitalizzare un segnale vocale con banda \([200 \; \text{Hz}, \; 3.4 \; \text{kHz}]\), si utilizza una frequenza
di campionamento \(f_c = 8 \; \text{kHz}\) e un quantizzatore da 256 livelli con 8 bpS (128 livelli con 7 bpS in Asia
e America). Eventualmente si suddividono i dati in pacchetti di 20 ms da 1.28 Mbit. Il bitrate risulta essere:
\[R_\text{voce} = f_c \cdot \log_2 L = 8 \; \text{kHz} \cdot 8 \; \text{bpS} = 64 \; \text{kbit/s} \qquad \text{packet size} = 20 \; \text{ms} \cdot R_\text{voce} = 1280 \; \text{bit} = 160 \; \text{byte}\]

\subsubsection*{Digitalizzazione di un segnale audio generico}
Per digitalizzare un segnale audio stereo generico (es. musica) con banda \([15 \; \text{Hz}, \; 20 \; \text{kHz}]\) si
utilizza una frequenza di campionamento di \(f_c = 44.1 \; \text{kHz}\) con 16 bpS, duplicato per i due canali. Il bitrate è di:
\[R_\text{musica} = 2 \cdot f_c \cdot \log_2 L = 2 \cdot 44.1 \; \text{kHz} \cdot 16 \; \text{bpS} = 1.411 \; \text{Mbit/s}\]

\subsubsection*{Memorizzazione dei segnali codificati con PCM nei CD}
Nei CD audio, l'informazione audio memorizzata è codificata tramite codifica PCM, e comprende solo l'87\% dei dati memorizzati.
Il restante 13\% è occupato dai codici di correzione d'errore (CRC). Il bitrate effettivo risulta essere:
\[R_\text{CD} = R_\text{musica} \; / \; 0.87 = 1.622 \; \text{Mbit/s}\]
Se si volesse calcolare la massima durata effettiva di un CD audio da 730 MB si esegue il seguente calcolo:
\[\text{durata} = \frac{\text{\# bit disponibili}}{\text{bit/s effettivi}} = \frac{730 \; \text{MB} \cdot 8 \cdot 10^6 \; \text{bit/MB}}{1.622 \cdot 10^6 \; \text{bit/s}} \approx 3600 \; \text{s} = 1 \; \text{h}\]

\subsection{Rappresentazione di immagini}
\subsubsection*{Rappresentazione di immagini B/N - campionamento spaziale e HDR}
La fase di acquisizione della luce di una scena reale tramite un sensore digitale (es. sensore CMOS) consiste in un
campionamento spaziale della luce che colpisce il sensore in ogni pixel di cui esso è costituito. Ogni pixel misura
l'intensità della luce che lo colpisce e la converte in un valore numerico, solitamente quantizzato su 8 bit da 0 a 255.
Il risultato è una matrice di valori numerici che rappresentano l'immagine acquisita in scala di grigi (B/N).

Esistono sensori digitali con gamma dinamica estesa e rappresentano le immagini in formato HDR in grado di supportare fino
a 10-12 bit per pixel. Tali tecnologie non vanno, però confuse con gli algoritmi che simulano l'HDR combinando più immagini
della stessa scena con esposizioni differenti.

\subsubsection*{Percezione del colore}
L'occhio umano percepisce la luce e i colori attraverso due tipi di cellule presenti sulla retina:
\begin{itemize}
	\item i bastoncelli, sono in numero maggiore e sono sensibili alla intensità luminosa
	\item i coni, sono in numero minore e sono sensibili a tre bande di frequenza della luce corrispondenti ai colori
	primari rosso (575 nm), verde (535 nm) e blu (445 nm)
\end{itemize}

\noindent
Si osserva, quindi, che è sufficiente memorizzare il contributo di ciascun colore primario per rappresentare un qualsiasi
colore percepito dall'occhio umano. Inoltre, essendo i coni in minoranza rispetto ai bastoncelli, l'occhio umano è più
sensibile alle variazioni di luminosità che alle variazioni di colore.

\subsubsection*{Rappresentazione di immagini a colori - RGB}
La rappresentazione di immagini a colori tramite il formato RGB consiste nel memorizzare il contributo di ciascun colore
primario (R, G, B) per ogni pixel dell'immagine. Ogni canale di colore è quantizzato su 8 bit con valori da 0 a 255, per
cui complessivamente ogni pixel è rappresentato da 24 bit e può codificare fino a \(2^{24} = 16\,777\,216\) colori diversi.

Anche per le immagini a colori esistono formati con gamma dinamica estesa (HDR) che supportano fino a 12 bit per canale
di colore, con 36 bit per pixel e \(2^{36} = 68\,719\,476\,736\) colori rappresentabili.

\subsubsection*{Rotazione dallo spazio RGB allo YCbCr}
Applicando un'opportuna rotazione degli assi, è possibile avere la diagonale principale \((0,0,0) \to (1,1,1)\) dello spazio
RGB come coordinata Y, con altre due coordinate ortogonali ad essa dette Cb e Cr. La coordinata Y è detta luminanza del colore,
infatti sulla diagonale principale si trovano tutte le tonalità di grigio, mentre le coordinate Cb e Cr sono il complemento
al blu e il complemento al rosso e rappresentano la componente cromatica detta anche crominanza del colore sul piano ortogonale
alla luminanza.

\subsubsection*{Sparsificazione del segnale YCbCr}
Un segnale viene detto sparso quando presenta alcuni campioni ``importanti'' e altri campioni ``meno importanti'' che
è possibile comprimere con apposite tecniche di sottocampionamento senza degradare eccessivamente il segnale rappresentato.

Il segnale YCbCr viene detto sparso in quanto l'occhio umano è più sensibile alla luminanza Y, rispetto al contributo
cromatico Cb e Cr. È possibile, quindi, applicare un processo di decimazione dei campioni Cb e Cr, riducendo il costo
di ogni pixel senza introdurre distorsioni percepibili dall'occhio umano.

\subsubsection*{Pipeline di conversione da RGB a YCbCr e viceversa}
\begin{figure}[htbp]
	\centering
	\begin{tikzpicture}[
		% Stile per i blocchi di elaborazione generici
		block/.style={
			rectangle,
			draw=blue!70!black,
			fill=blue!5,
			thick,
			align=center,
			rounded corners=3pt,
			font=\footnotesize
		},
		% Stile specifico per il canale di trasmissione
		channel/.style={
			rectangle,
			draw=orange!85!black,
			fill=orange!8,
			thick,
			align=center,
			rounded corners=3pt,
			font=\footnotesize
		},
		% Stile per le frecce di collegamento
		arrow/.style={
			-{Stealth[scale=1.2]},
			thick
		},
		% Stile per le etichette dei segnali
		signal/.style={
			align=center,
			font=\footnotesize
		},
		% Stile per le parentesi
		brace/.style={
			decorate,
			decoration={brace, amplitude=10pt, raise=6pt},
			thick
		},
		% Distanza standard tra i nodi
		%node distance=1.5cm and 1.8cm
	]
		% Catena di blocchi
		\node (in) at (0.5,0) [signal] {\textbf{Immagine}\\\textbf{RGB}};
		\node [block] at (4,0) (rgb2y) {RGB $\rightarrow$ YCbCr};
		\node [block] at (8,0) (decim) {Decimazione};
		\node [block] at (12,0) (write) {Scrittura file};
		\node [channel] at (12,-1) (storage) {Storage};
		\node [block] at (12,-2) (read) {Lettura file};
		\node [block] at (8,-2) (interp) {Interpolazione};
		\node [block] at (4,-2) (y2rgb) {YCbCr $\rightarrow$ RGB};
		\node [signal] at (0.5,-2) (out) {\textbf{Display}\\ (immagine RGB)};

		% Collegamenti ed etichette
		\draw [arrow] (in) -- (rgb2y);
		\draw [arrow] (rgb2y) -- node[signal, below] {Immagine\\YCbCr} (decim);
		\draw [arrow] (decim) -- node[signal, below] {Formato\\YUV 4:2:0} (write);
		\draw [arrow] (write) -- (storage);
		\draw [arrow] (storage) -- (read);
		\draw [arrow] (read) -- node[signal, above] {Formato\\YUV 4:2:0} (interp);
		\draw [arrow] (interp) -- node[signal, above] {Immagine\\YCbCr} (y2rgb);
		\draw [arrow] (y2rgb) -- (out);
	\end{tikzpicture}
\end{figure}

\noindent
Il processo di conversione da RGB a YCbCr e viceversa prevede i seguenti passaggi:
\begin{enumerate}
	\item conversione da RGB a YCbCr tramite una trasformazione lineare che ruota lo spazio dei colori, è implementabile
	come un semplice prodotto tra la matrice di trasformazione e il vettore RGB di ogni pixel dell'immagine
	\item decimazione dei campioni Cb e Cr, tramite il processo di chroma subsampling descritto in seguito
	\item memorizzazione in memoria o invio al destinatario del file compresso
	\item interpolazione per ricostruire i campioni Cb e Cr decimati, tramite un filtro di interpolazione
	\item conversione da YCbCr a RGB tramite la trasformazione lineare inversa
\end{enumerate}

\subsubsection*{Decimazione per sottocampionamento della crominanza (chroma subsampling)}
Il processo di decimazione dei campioni Cb e Cr viene effettuato tramite il formato YUV J:a:b. In pratica si suddivide
l'immagine in blocchi di \(4 \times 2\) pixel e i tre valori J, a e b indicano quanti campioni di ciascun canale vengono
memorizzati per ogni blocco. Di seguito alcuni formati di decimazione più diffusi, l'attuale standard universale è
il 4:2:0, utilizzato nella quasi totalità dei formati foto e video attualmente esistenti.
\begin{itemize}
	\item \textbf{4:4:4}: non viene effettuata alcuna decimazione
	\item \textbf{4:2:2}: per ogni riga si memorizzano solo 4 campioni di Y, 2 di Cb e 2 di Cr, riducendo la banda del
	colore del 50\% e il costo complessivo della rappresentazione del 33.3\%
	\item \textbf{4:2:0}: per ogni mezzo blocco di \(2 \times 2\) pixel si memorizzano solo 4 campioni di Y, 1 di Cb
	e 1 di Cr, riducendo la banda del colore del 75\% e il costo complessivo della rappresentazione del 50\%
	\item \textbf{4:1:1}: per ogni riga si memorizzano solo 4 campioni di Y, 1 di Cb e 1 di Cr, riducendo la banda del
	colore del 75\% e il costo complessivo della rappresentazione del 50\%
\end{itemize}

\noindent
In generale un'immagine HD di risoluzione \(1920 \times 1080\) rappresentata in formato RGB occupa approssimativamente
6.22 MBytes. Applicando la trasformazione da RGB a YCbCr e il sottocampionamento 4:2:0, il costo scende a 3.11 MBytes,
senza introdurre distorsioni percepibili dall'occhio umano.
\begin{align*}
	\text{size}(X_\text{RGB}) &= 1920 \times 1080 \; \text{pixel} \cdot 3 \; \text{canali} \cdot 8 \; \text{bit} \approx 49.77 \; \text{Mbit} = 6.22 \; \text{MByte} \\
	\text{size}(X_\text{YCbCr}) &= 1920 \times 1080 \; \text{pixel} \cdot 3 \; \text{canali} \cdot 8 \; \text{bit} / 2 \approx 24.88 \; \text{Mbit} = 3.11 \; \text{MByte}
\end{align*}

\subsection{Rappresentazione di video}
\subsubsection*{Percezione del movimento}
Il fenomeno della persistenza retinale consiste nel fatto che un'immagine osservata permane per un breve periodo di tempo
sulla retina dell'occhio umano. Questo permette di campionare nel tempo un segnale video e di visualizzare le immagini
campionate in rapida successione, dando l'illusione di un movimento continuo. Le immagini campionate sono dette fotogrammi
o frame e solitamente vengono visualizzate a 24, 30, 50 o 60 frame al secondo (fps).

\subsubsection*{Formati video non compressi}
Per un video non compresso, il bitrate è estremamente elevato e si richiederanno altre tecniche di compressione per poterlo
trasmettere efficacemente in rete. Ad esempio un video HD di risoluzione \(1920 \times 1080\) con 50 fps con frame codificati
in formato YCbCr 4:2:0 ha un bitrate di circa 155 MBytes/s.
\[R_\text{YCbCr} = \text{frame-size}_\text{\;YCbCr} \cdot 50 \; \text{fps} = 3.11 \; \text{MBytes} \cdot 50 \; \text{fps} \approx 1.244 \; \text{Gbit/s} = 155.5 \; \text{MBytes/s}\]
